\section{Related Work}
\label{sec:related-work}

\paragraph{Real-to-sim reconstruction.}
RialTo scans a scene into a digital twin for policy robustification
\citep{torne2024rialto}. Digital cousins build interactive lookalikes from a
single RGB image \citep{dai2024cousins}. SIMPLER matches real robot setups so
that simulated evaluation tracks real evaluation \citep{li2024simpler}.
URDFormer predicts articulated URDF scenes directly from real photographs
\citep{chen2024urdformer}. Lab2Sim shares the goal of an interactive twin but
takes a different input: a protocol and the manufacturer's own evidence for
the named instrument, rather than a scan or photograph of a workspace that
already exists.

\paragraph{Generative scene agents.}
GenSim writes simulation task code from language \citep{wang2024gensim}.
RoboGen proposes tasks, scenes, and supervision in a propose--generate--learn
loop \citep{wang2024robogen}. Holodeck lays out 3D rooms from a prompt using
Objaverse meshes \citep{yang2024holodeck}. Image-conditioned 3D generators such
as Hunyuan3D propose textured geometry from a single view
\citep{zhao2025hunyuan3d}. Lab2Sim calls generators of this kind as tools
inside its Blender skill. What it adds is the layer these agents leave
implicit: a written specification of the instrument's joints and contact
surfaces, and an admission step that decides whether the generated USD may
enter the robot's laboratory.
